%% notes-tex/010-additive-baselines/sections/5-comparability.tex

\section{Is the comparison fair\secstatus{todo}}
\label{sec:comparability}

Six things have to match before a baseline number can sit in the same table as a
transformer number. Table~\ref{tab:comparability} names them and gives the
verdict on each; the subsections that follow give the evidence.

%% SOURCE: each row is established in the subsection it names; the
%% normalization constants are read from the runs' logged transform config.
\begin{table}[htbp]
\centering
\begin{tabular}{llp{0.42\textwidth}}
\hline
\# & Verdict & What has to match \\
\hline
1 & matches & \textbf{Records and split.} All models use the same 376{,}732
    records and the same \texttt{index\_seed\_42.json} partition. \\
2 & matches & \textbf{Label.} All predict \texttt{gene\_interaction}, the
    published adjusted $\tau$ of Eq.~\ref{eq:tau}. \\
3 & matches & \textbf{Reported split.} Test is reported for every row, and
    validation is what every row selected on. \\
4 & matches & \textbf{Metric definition and denominator.} Pearson, Spearman and
    mean squared error over the full 37{,}673 test records in raw $\tau$ units
    for every row. \\
5 & matches & \textbf{Information available.} The identity of the three
    perturbed genes, and nothing else, for every model
    (Section~\ref{sec:sees}). \\
6 & \textbf{does not} & \textbf{Label normalization.} The transformer trains on
    z-scored labels whose mean and standard deviation were computed once, before
    the split, from all 376{,}732 records; the baselines use the training mean
    only. Pearson is invariant to the two constants, so this cannot move the
    reported metric (Section~\ref{sec:normleak}). \\
\hline
\end{tabular}
\caption{What has to match for the comparison to be fair. Row 6 is the one
mismatch. It is a hygiene defect in the 010 runs rather than a source of
advantage, and the 025 arms from the disjoint run onward fit on the training
split.}
\label{tab:comparability}
\end{table}

The loss functions also differ between the baselines and the transformer, but
that is a difference in what each model was trained to do rather than in how the
comparison is scored, so it is not one of the six and is treated separately
below.

\subsection{Which split the reported numbers come from\secstatus{todo}}

Both validation and test numbers are available for the transformer, and the
table in Section~\ref{sec:results} reports \textbf{test}. The distinction
matters because the three transformer checkpoints are \emph{best-Pearson}
checkpoints, selected on validation across roughly 25 to 64 validation epochs.
A validation Pearson chosen that way is the maximum of many draws and is
upward-biased as an estimate of capability; its test Pearson is not, because
test was not consulted during selection.

The baselines follow the same discipline. B1, B2 and B4 select the ridge
penalty on validation and report test. B5 early-stops on validation and reports
test. So the test column compares like with like, and the validation column is
biased in the same direction for every row.

\subsection{What the transformer metrics are\secstatus{todo}}

The metrics come from the three checkpoints' re-evaluation runs under
\texttt{\$DATA\_ROOT/wandb-experiments}, read from
\texttt{files/wandb-summary.json} rather than retyped. Two families of key exist
and they are not interchangeable.

\begin{itemize}
\item \texttt{val/gene\_interaction/\{MSE,RMSE,Pearson\}} and the \texttt{test/}
  equivalents are torchmetrics objects updated once per batch and computed at
  epoch end, so they cover the \emph{full} 37{,}673 records of each split. They
  are computed after the inverse transform, in raw $\tau$ units.
\item \texttt{val\_sample/Spearman\_target\_0} and its test counterpart come
  from a separate path that collects samples for plotting and is capped by
  \texttt{plot\_sample\_ceiling}. In the evaluation configuration that ceiling
  is $10^{6}$, above the split size, so no truncation fires and the Spearman is
  also over the full split. In the \emph{training} configuration the same
  ceiling is 10{,}000, so a Spearman quoted from a training run would be over a
  random 10{,}000 of 37{,}673. The numbers here are from the evaluation runs.
\item \texttt{val/transformed/gene\_interaction/*} are the same metrics in
  standardized units and are not used here.
\end{itemize}

The baselines compute Pearson, Spearman and mean squared error over the same
full splits in raw $\tau$ units, so the units and the denominators agree.

\subsection{What each model sees\secstatus{todo}}

Established in Section~\ref{sec:sees}: the perturbed gene set, and nothing else,
for every model in the comparison. No baseline is handicapped by being denied a
feature the transformer had, and the transformer is not credited with
information it never received.

\subsection{Where the comparison is not clean\secstatus{todo}}

\notesec{Label normalization: what it is, when it happens, and what it can move}
\label{sec:normleak}
The transformer is trained on the z-scored label $(y - \mu) / \sigma$, the
\texttt{standard} strategy of \texttt{COOLabelNormalizationTransform}. The two
constants are computed once, at start-up, from \texttt{dataset.label\_df}, and
the transform is built before the data module splits the records, so in 010 the
population is every record in the build. The logged constants confirm it:
$\mu = -0.008024324$ and $\sigma = 0.063263549$ are the all-record statistics to
nine digits, where the training split alone gives $-0.007688739$ and
$0.063186255$. The inverse transform is applied to every prediction before a
metric is computed, so all reported numbers are in raw $\tau$ units.

What the choice of population can and cannot do. Pearson and Spearman are
invariant to an affine map of the labels, so no choice of $\mu$ and $\sigma$
changes the correlation a given model would reach; the two validation and test
scalars that reach training this way carry no information a model could turn
into held-out correlation. What the constants do set is the scale of the point
loss relative to the two other terms of Eq.~\ref{eq:loss}: a different $\sigma$
multiplies the standardized mean squared error by $(\sigma_{\mathrm{all}} /
\sigma_{\mathrm{train}})^2 = 1.0024$, a quarter of a percent, and a different
$\mu$ is a constant offset the readout bias absorbs. Against a graph term that is
99.99 percent of the objective (Section~\ref{sec:graphs}) that is nothing. So
row 6 of Table~\ref{tab:comparability} is a discipline defect, held-out labels
touching a training-time constant, and not a mechanism for the transformer's
margin.

It is corrected from the disjoint run onward. The 025 trainer now names the
fit population, \texttt{transforms.fit\_on}, and the query-pair-disjoint arm
fits on its 301{,}236 training records, where the constants are
$\mu = -0.007843739$ and $\sigma = 0.063776418$; the 025 replication arm keeps
\texttt{fit\_on: subset} because reproducing 010 requires 010's constants.
Both sets are written by
\texttt{experiments/025-solid-growth/scripts/label\_normalization\_constants.py}.

\notesec{The loss functions differ, deliberately}
The baselines minimize squared error. The transformer minimizes
%
\begin{equation}
  \mathcal{L} \;=\; \lambda_p \mathcal{L}_{\mathrm{MSE}}
  \;+\; \lambda_d\, S_{\epsilon}\!\left(
    \tfrac{1}{B}\textstyle\sum_b \delta_{\hat{y}_b},\;
    \tfrac{1}{B}\textstyle\sum_b \delta_{y_b}\right)
  \;+\; \lambda_g \mathcal{L}_{\mathrm{graph}},
  \label{eq:loss}
\end{equation}
%
with $\lambda_p = 1$, $\lambda_d = 0.1$, $\lambda_g = 1$. The middle term is a
debiased Sinkhorn divergence with $p = 2$ and blur $0.05$ between the empirical
measure of the batch's predictions and that of the batch's targets, computed
unpaired, which shapes the marginal distribution of the predictions rather than
their per-record accuracy. The third term is the graph penalty of
Section~\ref{sec:graphs}.

This is a difference in training objective, not in the evaluation, and it is
part of what the transformer is. It is noted because a distributional term can
move Pearson and mean squared error in opposite directions, so the two metrics
should be read together rather than either alone.

\notesec{Paired residuals had to be recomputed}
The original evaluation runs wrote per-record prediction files to the cluster
checkout, and those files are not on this machine. Recomputing them through the
stock evaluation script fails on the frozen build, which no longer validates
(Section~\ref{sec:modernize}). Section~\ref{sec:verify} gives the route around
that, and Section~\ref{sec:modernize} gives the migration that removes the
obstacle for future work.

\subsection{Model size\secstatus{todo}}

%% SOURCE: parameter counts read from
%% torchcell/models/equivariant_cell_graph_transformer.py at the lines given
\begin{table}[htbp]
\centering
\begin{tabular}{lrl}
\hline
Component & Parameters & Role \\
\hline
Gene embedding table & 1{,}189{,}260 & the only place gene identity lives \\
Encoder, 8 layers & 3{,}129{,}120 & strain-independent, see Sec.~\ref{sec:encoder} \\
Perturbation transform & 391{,}140 & the strain-dependent step \\
Readout head & 65{,}161 & pooling and MLP \\
CLS token & 180 & constant \\
\hline
Total & 4{,}774{,}861 & \\
\hline
B1 additive ridge & 4{,}353 & one scalar per gene, plus intercept \\
B5 embedding MLP & 623{,}745 & 128 per gene, plus a two-layer network \\
\hline
\end{tabular}
\caption{Parameter budget. Two thirds of the transformer sits in the encoder,
which does no per-sample work, and a quarter in the embedding table. The
strain-dependent computation is 456{,}301 parameters, under a tenth of the
model.}
\label{tab:params}
\end{table}
